% Nonlinear Control Theory Template
% Topics: Lyapunov stability, phase plane analysis, feedback linearization, sliding mode control, backstepping
% Style: Advanced control systems analysis with computational verification

\documentclass[a4paper, 11pt]{article}
\usepackage[utf8]{inputenc}
\usepackage[T1]{fontenc}
\usepackage{amsmath, amssymb}
\usepackage{graphicx}
\usepackage{siunitx}
\usepackage{booktabs}
\usepackage{subcaption}
\usepackage[makestderr]{pythontex}
\usepackage{hyperref}

% Theorem environments
\newtheorem{definition}{Definition}[section]
\newtheorem{theorem}{Theorem}[section]
\newtheorem{lemma}{Lemma}[section]
\newtheorem{example}{Example}[section]
\newtheorem{remark}{Remark}[section]

\title{Nonlinear Control Systems: Stability Analysis and Advanced Control Design}
\author{Control Systems Laboratory}
\date{\today}

\begin{document}
\maketitle

\begin{abstract}
This report presents a comprehensive analysis of nonlinear control systems using Lyapunov
stability theory, phase plane methods, and advanced nonlinear control design techniques.
We examine the stability of equilibrium points for representative nonlinear systems,
demonstrate feedback linearization for affine nonlinear systems, design sliding mode
controllers with chattering reduction, and apply backstepping to cascade systems.
Computational analysis verifies stability margins, control performance, and robustness
properties through simulation of Van der Pol oscillator, inverted pendulum, and nonlinear
mass-spring-damper systems.
\end{abstract}

\section{Introduction}

Nonlinear control theory addresses systems that cannot be adequately described by linear
approximations, particularly when operating over wide ranges or exhibiting inherently
nonlinear phenomena such as limit cycles, bifurcations, and chaotic behavior. Unlike
linear control, where superposition and frequency-domain methods dominate, nonlinear
control relies on state-space methods, Lyapunov theory, and differential geometric tools.

\begin{definition}[Nonlinear Dynamical System]
A continuous-time nonlinear dynamical system is described by:
\begin{equation}
\dot{\mathbf{x}} = \mathbf{f}(\mathbf{x}, \mathbf{u}, t)
\end{equation}
where $\mathbf{x} \in \mathbb{R}^n$ is the state vector, $\mathbf{u} \in \mathbb{R}^m$ is the control input,
and $\mathbf{f}: \mathbb{R}^n \times \mathbb{R}^m \times \mathbb{R} \rightarrow \mathbb{R}^n$ is generally nonlinear.
\end{definition}

\section{Lyapunov Stability Theory}

\subsection{Direct Method of Lyapunov}

\begin{theorem}[Lyapunov Stability]
Consider the autonomous system $\dot{\mathbf{x}} = \mathbf{f}(\mathbf{x})$ with equilibrium point
$\mathbf{x}_e$ (i.e., $\mathbf{f}(\mathbf{x}_e) = \mathbf{0}$). If there exists a continuously differentiable
function $V: D \rightarrow \mathbb{R}$ where $D$ contains $\mathbf{x}_e$, such that:
\begin{enumerate}
\item $V(\mathbf{x}_e) = 0$ and $V(\mathbf{x}) > 0$ for all $\mathbf{x} \neq \mathbf{x}_e$ (positive definite)
\item $\dot{V}(\mathbf{x}) = \nabla V \cdot \mathbf{f}(\mathbf{x}) \leq 0$ (negative semidefinite)
\end{enumerate}
then $\mathbf{x}_e$ is stable. If additionally $\dot{V}(\mathbf{x}) < 0$ for all $\mathbf{x} \neq \mathbf{x}_e$
(negative definite), then $\mathbf{x}_e$ is asymptotically stable.
\end{theorem}

\begin{definition}[Control Lyapunov Function (CLF)]
For a control system $\dot{\mathbf{x}} = \mathbf{f}(\mathbf{x}, \mathbf{u})$, a function $V(\mathbf{x})$ is
a control Lyapunov function if there exists a control law $\mathbf{u} = \mathbf{k}(\mathbf{x})$ such that
$\dot{V}(\mathbf{x}) < 0$ for all $\mathbf{x} \neq \mathbf{0}$.
\end{definition}

\subsection{Phase Plane Analysis}

\begin{definition}[Equilibrium Point Classification]
For a two-dimensional system $\dot{\mathbf{x}} = \mathbf{f}(\mathbf{x})$, equilibrium points are
classified by linearization eigenvalues $\lambda_1, \lambda_2$:
\begin{itemize}
\item \textbf{Stable node}: Both $\lambda_i < 0$ (real)
\item \textbf{Unstable node}: Both $\lambda_i > 0$ (real)
\item \textbf{Saddle point}: $\lambda_1 < 0 < \lambda_2$ (real)
\item \textbf{Stable focus}: $\text{Re}(\lambda_i) < 0$ (complex)
\item \textbf{Center}: $\text{Re}(\lambda_i) = 0$ (pure imaginary)
\end{itemize}
\end{definition}

\begin{example}[Van der Pol Oscillator]
The Van der Pol equation models self-excited oscillations:
\begin{equation}
\ddot{x} - \mu(1 - x^2)\dot{x} + x = 0
\end{equation}
Rewritten as $\dot{x}_1 = x_2$, $\dot{x}_2 = \mu(1 - x_1^2)x_2 - x_1$, this exhibits a stable
limit cycle for $\mu > 0$.
\end{example}

\section{Feedback Linearization}

\begin{definition}[Affine Nonlinear System]
A single-input affine nonlinear system has the form:
\begin{equation}
\dot{\mathbf{x}} = \mathbf{f}(\mathbf{x}) + \mathbf{g}(\mathbf{x})u
\end{equation}
where $\mathbf{f}, \mathbf{g}: \mathbb{R}^n \rightarrow \mathbb{R}^n$ are smooth vector fields.
\end{definition}

\begin{theorem}[Input-Output Linearization]
For output $y = h(\mathbf{x})$, compute the relative degree $r$ as the smallest integer such that:
\begin{equation}
\frac{\partial (L_{\mathbf{f}}^{r-1} h)}{\partial \mathbf{x}} \mathbf{g}(\mathbf{x}) \neq 0
\end{equation}
where $L_{\mathbf{f}} h = \nabla h \cdot \mathbf{f}$ denotes the Lie derivative. The control law:
\begin{equation}
u = \frac{1}{L_{\mathbf{g}} L_{\mathbf{f}}^{r-1} h} \left( v - L_{\mathbf{f}}^r h \right)
\end{equation}
linearizes the input-output map to $y^{(r)} = v$.
\end{theorem}

\section{Sliding Mode Control}

\begin{definition}[Sliding Surface]
A sliding surface $s(\mathbf{x}) = 0$ defines a manifold in state space where the controlled
system exhibits desired dynamics. The control objective is to drive $s(\mathbf{x}) \rightarrow 0$
in finite time and maintain it thereafter.
\end{definition}

\begin{theorem}[Reaching Condition]
The sliding condition $s \dot{s} < 0$ ensures trajectories reach the sliding surface in finite time.
A common choice is:
\begin{equation}
\dot{s} = -k \, \text{sign}(s), \quad k > 0
\end{equation}
yielding reaching time $t_r = |s(0)| / k$.
\end{theorem}

\begin{remark}[Chattering Reduction]
Discontinuous control causes chattering. Continuous approximations like:
\begin{equation}
u = -k \frac{s}{\epsilon + |s|}
\end{equation}
reduce chattering while preserving stability with practical precision $\epsilon$.
\end{remark}

\section{Backstepping}

\begin{definition}[Strict-Feedback Form]
A system is in strict-feedback form if:
\begin{align}
\dot{x}_1 &= x_2 + \phi_1(x_1) \\
\dot{x}_2 &= x_3 + \phi_2(x_1, x_2) \\
&\vdots \\
\dot{x}_n &= u + \phi_n(x_1, \ldots, x_n)
\end{align}
\end{definition}

\begin{theorem}[Recursive Backstepping]
Stabilizing controls are constructed recursively. At step $i$, treat $x_{i+1}$ as virtual control,
design $\alpha_i(x_1, \ldots, x_i)$ to stabilize subsystem $\{x_1, \ldots, x_i\}$, then augment
the CLF to account for error $z_i = x_i - \alpha_{i-1}$.
\end{theorem}

\section{Computational Analysis}

\begin{pycode}
import numpy as np
import matplotlib.pyplot as plt
from scipy.integrate import solve_ivp
from matplotlib import cm
from matplotlib.patches import FancyArrowPatch

np.random.seed(42)

# ==============================================================================
# 1. Lyapunov Stability for Nonlinear System
# ==============================================================================

def nonlinear_system(t, x, control_on=False):
    """Nonlinear system: x1' = -x1 + x2, x2' = -x1*x2 - x2^3 + u"""
    x1, x2 = x
    if control_on:
        # Lyapunov-based control: u = -k*(∂V/∂x2)
        V_grad_x2 = 2*x2
        u = -2.0 * V_grad_x2
    else:
        u = 0
    dx1dt = -x1 + x2
    dx2dt = -x1*x2 - x2**3 + u
    return [dx1dt, dx2dt]

def lyapunov_V(x1, x2):
    """Candidate Lyapunov function V = x1^2 + x2^2"""
    return x1**2 + x2**2

def lyapunov_Vdot(x1, x2, control_on=False):
    """Time derivative of V along trajectories"""
    dx = nonlinear_system(0, [x1, x2], control_on)
    return 2*x1*dx[0] + 2*x2*dx[1]

# ==============================================================================
# 2. Van der Pol Oscillator Phase Portrait
# ==============================================================================

def van_der_pol(t, x, mu):
    """Van der Pol oscillator"""
    x1, x2 = x
    dx1dt = x2
    dx2dt = mu*(1 - x1**2)*x2 - x1
    return [dx1dt, dx2dt]

# ==============================================================================
# 3. Feedback Linearization: Inverted Pendulum
# ==============================================================================

def pendulum_dynamics(t, x, control_law='none'):
    """Inverted pendulum: θ'' = sin(θ) + u/m*L^2"""
    theta, theta_dot = x
    m, L, g = 1.0, 1.0, 9.81

    if control_law == 'feedback_linearization':
        # Exact linearization: u = -m*L^2*(sin(θ) + kd*θ' + kp*θ)
        kp, kd = 20.0, 8.0
        u = -m*L**2*(np.sin(theta) + kd*theta_dot + kp*theta)
        ddtheta = np.sin(theta) + u/(m*L**2)
    else:
        # Uncontrolled
        ddtheta = np.sin(theta)

    return [theta_dot, ddtheta]

# ==============================================================================
# 4. Sliding Mode Control: Double Integrator
# ==============================================================================

def double_integrator_smc(t, x, epsilon=0.1):
    """Double integrator with sliding mode control"""
    x1, x2 = x
    # Sliding surface: s = x2 + λ*x1
    lam = 2.0
    s = x2 + lam*x1

    # Control law: u = -k*sat(s/ε)
    k = 5.0
    u = -k * np.tanh(s/epsilon)  # Smooth approximation

    dx1dt = x2
    dx2dt = u
    return [dx1dt, dx2dt, s]

# ==============================================================================
# 5. Backstepping: Nonlinear Mass-Spring-Damper
# ==============================================================================

def backstepping_system(t, x):
    """Strict-feedback system with backstepping control"""
    x1, x2 = x

    # System: x1' = x2, x2' = -0.5*x2 - sin(x1) + u
    # Step 1: Virtual control α1 = -c1*x1
    c1 = 2.0
    alpha1 = -c1 * x1

    # Step 2: Error z2 = x2 - α1
    z2 = x2 - alpha1

    # Step 3: Actual control u
    c2 = 3.0
    u = -np.sin(x1) - 0.5*x2 - z2 - c2*z2 + c1*x2

    dx1dt = x2
    dx2dt = -0.5*x2 - np.sin(x1) + u
    return [dx1dt, dx2dt]

# ==============================================================================
# SIMULATIONS
# ==============================================================================

t_span = (0, 20)
t_eval = np.linspace(0, 20, 1000)

# 1. Lyapunov stability comparison
x0_lyap = [2.0, 1.5]
sol_uncontrolled = solve_ivp(lambda t, x: nonlinear_system(t, x, False),
                              t_span, x0_lyap, t_eval=t_eval, max_step=0.01)
sol_controlled = solve_ivp(lambda t, x: nonlinear_system(t, x, True),
                           t_span, x0_lyap, t_eval=t_eval, max_step=0.01)

# Compute Lyapunov function values
V_uncontrolled = lyapunov_V(sol_uncontrolled.y[0], sol_uncontrolled.y[1])
V_controlled = lyapunov_V(sol_controlled.y[0], sol_controlled.y[1])

# 2. Van der Pol with limit cycle
mu_vdp = 2.0
x0_vdp = [0.5, 0.0]
sol_vdp = solve_ivp(lambda t, x: van_der_pol(t, x, mu_vdp),
                    (0, 30), x0_vdp, max_step=0.01)

# 3. Pendulum feedback linearization
x0_pend = [np.pi - 0.3, 0.1]  # Near inverted position
sol_pend_unc = solve_ivp(lambda t, x: pendulum_dynamics(t, x, 'none'),
                         (0, 5), x0_pend, t_eval=np.linspace(0, 5, 500), max_step=0.01)
sol_pend_fl = solve_ivp(lambda t, x: pendulum_dynamics(t, x, 'feedback_linearization'),
                        (0, 5), x0_pend, t_eval=np.linspace(0, 5, 500), max_step=0.01)

# 4. Sliding mode control
x0_smc = [3.0, 1.0]
t_smc = np.linspace(0, 10, 1000)
sol_smc = solve_ivp(lambda t, x: double_integrator_smc(t, x)[:2],
                    (0, 10), x0_smc, t_eval=t_smc, max_step=0.01)

# Compute sliding surface
sliding_surface = sol_smc.y[1] + 2.0*sol_smc.y[0]

# 5. Backstepping
x0_back = [1.5, -0.5]
sol_back = solve_ivp(backstepping_system, (0, 10), x0_back,
                     t_eval=np.linspace(0, 10, 500), max_step=0.01)

# ==============================================================================
# PLOTTING
# ==============================================================================

fig = plt.figure(figsize=(16, 14))

# Plot 1: Lyapunov function decrease
ax1 = fig.add_subplot(3, 3, 1)
ax1.plot(sol_uncontrolled.t, V_uncontrolled, 'r-', linewidth=2, label='Uncontrolled')
ax1.plot(sol_controlled.t, V_controlled, 'b-', linewidth=2, label='Lyapunov Control')
ax1.set_xlabel('Time (s)', fontsize=10)
ax1.set_ylabel('$V(\\mathbf{x}) = x_1^2 + x_2^2$', fontsize=10)
ax1.set_title('Lyapunov Function Evolution', fontsize=11, fontweight='bold')
ax1.legend(fontsize=9)
ax1.grid(True, alpha=0.3)
ax1.set_ylim([0, max(V_uncontrolled.max(), V_controlled.max())*1.1])

# Plot 2: Phase portrait with Lyapunov contours
ax2 = fig.add_subplot(3, 3, 2)
x1_range = np.linspace(-3, 3, 100)
x2_range = np.linspace(-3, 3, 100)
X1, X2 = np.meshgrid(x1_range, x2_range)
V_field = lyapunov_V(X1, X2)
Vdot_field = np.zeros_like(V_field)
for i in range(len(x1_range)):
    for j in range(len(x2_range)):
        Vdot_field[j, i] = lyapunov_Vdot(X1[j, i], X2[j, i], True)

contour = ax2.contour(X1, X2, V_field, levels=10, colors='gray', alpha=0.4, linewidths=0.8)
ax2.clabel(contour, inline=True, fontsize=7)
ax2.plot(sol_controlled.y[0], sol_controlled.y[1], 'b-', linewidth=2, label='Controlled')
ax2.plot(sol_uncontrolled.y[0], sol_uncontrolled.y[1], 'r--', linewidth=1.5, label='Uncontrolled')
ax2.scatter([0], [0], s=100, c='green', marker='*', edgecolor='black', zorder=5, label='Equilibrium')
ax2.set_xlabel('$x_1$', fontsize=10)
ax2.set_ylabel('$x_2$', fontsize=10)
ax2.set_title('Phase Portrait with Lyapunov Contours', fontsize=11, fontweight='bold')
ax2.legend(fontsize=8)
ax2.grid(True, alpha=0.3)
ax2.set_xlim([-3, 3])
ax2.set_ylim([-3, 3])

# Plot 3: Lyapunov derivative field
ax3 = fig.add_subplot(3, 3, 3)
cs = ax3.contourf(X1, X2, Vdot_field, levels=20, cmap='RdBu_r', vmin=-10, vmax=10)
cbar = plt.colorbar(cs, ax=ax3)
cbar.set_label('$\\dot{V}(\\mathbf{x})$', fontsize=9)
ax3.contour(X1, X2, Vdot_field, levels=[0], colors='black', linewidths=2)
ax3.set_xlabel('$x_1$', fontsize=10)
ax3.set_ylabel('$x_2$', fontsize=10)
ax3.set_title('$\\dot{V}$ Field (Negative = Stable)', fontsize=11, fontweight='bold')

# Plot 4: Van der Pol phase portrait
ax4 = fig.add_subplot(3, 3, 4)
ax4.plot(sol_vdp.y[0], sol_vdp.y[1], 'b-', linewidth=1.5)
ax4.scatter([sol_vdp.y[0][0]], [sol_vdp.y[1][0]], s=80, c='green', marker='o',
           edgecolor='black', zorder=5, label='Initial')
ax4.scatter([sol_vdp.y[0][-1]], [sol_vdp.y[1][-1]], s=80, c='red', marker='s',
           edgecolor='black', zorder=5, label='Final')
ax4.scatter([0], [0], s=100, c='orange', marker='*', edgecolor='black', zorder=5,
           label='Equilibrium')
ax4.set_xlabel('$x_1$', fontsize=10)
ax4.set_ylabel('$x_2 = \\dot{x}_1$', fontsize=10)
ax4.set_title(f'Van der Pol Limit Cycle ($\\mu = {mu_vdp}$)', fontsize=11, fontweight='bold')
ax4.legend(fontsize=8)
ax4.grid(True, alpha=0.3)

# Plot 5: Van der Pol time series
ax5 = fig.add_subplot(3, 3, 5)
ax5.plot(sol_vdp.t, sol_vdp.y[0], 'b-', linewidth=1.5, label='$x_1$')
ax5.plot(sol_vdp.t, sol_vdp.y[1], 'r--', linewidth=1.5, label='$x_2$')
ax5.set_xlabel('Time (s)', fontsize=10)
ax5.set_ylabel('State', fontsize=10)
ax5.set_title('Van der Pol Time Response', fontsize=11, fontweight='bold')
ax5.legend(fontsize=9)
ax5.grid(True, alpha=0.3)

# Plot 6: Pendulum feedback linearization
ax6 = fig.add_subplot(3, 3, 6)
ax6.plot(sol_pend_unc.t, np.rad2deg(sol_pend_unc.y[0]), 'r--',
        linewidth=2, label='Uncontrolled')
ax6.plot(sol_pend_fl.t, np.rad2deg(sol_pend_fl.y[0]), 'b-',
        linewidth=2, label='Feedback Linearization')
ax6.axhline(y=180, color='green', linestyle=':', linewidth=1.5, label='Target')
ax6.set_xlabel('Time (s)', fontsize=10)
ax6.set_ylabel('$\\theta$ (degrees)', fontsize=10)
ax6.set_title('Inverted Pendulum Stabilization', fontsize=11, fontweight='bold')
ax6.legend(fontsize=9)
ax6.grid(True, alpha=0.3)

# Plot 7: Sliding mode control - phase plane
ax7 = fig.add_subplot(3, 3, 7)
ax7.plot(sol_smc.y[0], sol_smc.y[1], 'b-', linewidth=2, label='Trajectory')
# Plot sliding surface
x1_slide = np.linspace(-3.5, 3.5, 100)
x2_slide = -2.0 * x1_slide
ax7.plot(x1_slide, x2_slide, 'r--', linewidth=2, label='Sliding Surface')
ax7.scatter([sol_smc.y[0][0]], [sol_smc.y[1][0]], s=80, c='green',
           marker='o', edgecolor='black', zorder=5)
ax7.scatter([0], [0], s=100, c='orange', marker='*', edgecolor='black', zorder=5)
ax7.set_xlabel('$x_1$', fontsize=10)
ax7.set_ylabel('$x_2$', fontsize=10)
ax7.set_title('Sliding Mode Control Phase Portrait', fontsize=11, fontweight='bold')
ax7.legend(fontsize=9)
ax7.grid(True, alpha=0.3)
ax7.set_xlim([-3.5, 3.5])
ax7.set_ylim([-3, 3])

# Plot 8: Sliding surface evolution
ax8 = fig.add_subplot(3, 3, 8)
ax8.plot(t_smc, sliding_surface, 'b-', linewidth=2)
ax8.axhline(y=0, color='r', linestyle='--', linewidth=1.5)
ax8.fill_between(t_smc, -0.1, 0.1, alpha=0.2, color='green',
                label='Sliding Region')
ax8.set_xlabel('Time (s)', fontsize=10)
ax8.set_ylabel('$s = x_2 + 2x_1$', fontsize=10)
ax8.set_title('Sliding Surface Convergence', fontsize=11, fontweight='bold')
ax8.legend(fontsize=9)
ax8.grid(True, alpha=0.3)

# Plot 9: Backstepping control
ax9 = fig.add_subplot(3, 3, 9)
ax9.plot(sol_back.t, sol_back.y[0], 'b-', linewidth=2, label='$x_1$')
ax9.plot(sol_back.t, sol_back.y[1], 'r--', linewidth=2, label='$x_2$')
ax9.axhline(y=0, color='black', linestyle=':', linewidth=1)
ax9.set_xlabel('Time (s)', fontsize=10)
ax9.set_ylabel('State', fontsize=10)
ax9.set_title('Backstepping Control Response', fontsize=11, fontweight='bold')
ax9.legend(fontsize=9)
ax9.grid(True, alpha=0.3)

plt.tight_layout()
plt.savefig('nonlinear_control_analysis.pdf', dpi=150, bbox_inches='tight')
plt.close()

# ==============================================================================
# Performance Metrics
# ==============================================================================

# Lyapunov control performance
settling_time_lyap = sol_controlled.t[np.where(V_controlled < 0.01)[0][0]] if any(V_controlled < 0.01) else sol_controlled.t[-1]
final_error_lyap = np.sqrt(V_controlled[-1])

# Pendulum stabilization error
theta_error_fl = np.abs(sol_pend_fl.y[0] - np.pi)
settling_time_pend = sol_pend_fl.t[np.where(theta_error_fl < 0.05)[0][0]] if any(theta_error_fl < 0.05) else sol_pend_fl.t[-1]

# Sliding mode reaching time
reaching_time_smc = t_smc[np.where(np.abs(sliding_surface) < 0.1)[0][0]] if any(np.abs(sliding_surface) < 0.1) else t_smc[-1]

# Backstepping final error
final_error_back = np.sqrt(sol_back.y[0][-1]**2 + sol_back.y[1][-1]**2)

\end{pycode}

\begin{figure}[htbp]
\centering
\includegraphics[width=\textwidth]{nonlinear_control_analysis.pdf}
\caption{Comprehensive nonlinear control analysis: (a) Lyapunov function $V(\mathbf{x}) = x_1^2 + x_2^2$
decreasing to zero under Lyapunov-based control, demonstrating asymptotic stability with exponential
convergence rate; (b) Phase portrait showing controlled trajectory (blue) converging to equilibrium
along level curves of the Lyapunov function, while uncontrolled trajectory (red dashed) diverges;
(c) Contour plot of $\dot{V}(\mathbf{x})$ showing negative semidefinite region (blue) where stability
is guaranteed; (d) Van der Pol oscillator exhibiting stable limit cycle for $\mu = 2.0$, demonstrating
self-excited oscillations characteristic of nonlinear systems; (e) Time-domain response showing periodic
oscillation settling into limit cycle; (f) Inverted pendulum stabilization at $\theta = 180°$ via
feedback linearization, achieving exact cancellation of nonlinear terms and exponential convergence
with control gains $k_p = 20$, $k_d = 8$; (g) Sliding mode control phase portrait showing trajectory
reaching sliding surface $s = x_2 + 2x_1 = 0$ in finite time then sliding to origin; (h) Sliding
surface variable converging to zero with smooth saturation boundary layer to eliminate chattering;
(i) Backstepping control for strict-feedback system achieving asymptotic stabilization through
recursive Lyapunov design with virtual control errors driving both states to zero.}
\label{fig:nonlinear_control}
\end{figure}

\section{Results}

\subsection{Stability Margins and Control Performance}

\begin{pycode}
print(r"\begin{table}[htbp]")
print(r"\centering")
print(r"\caption{Control Performance Metrics}")
print(r"\begin{tabular}{lcccc}")
print(r"\toprule")
print(r"Control Method & Settling Time (s) & Final Error & Control Gain & Stability Type \\")
print(r"\midrule")

print(f"Lyapunov-based & {settling_time_lyap:.2f} & {final_error_lyap:.4f} & $k = 2.0$ & Asymptotic \\\\")
print(f"Feedback Linearization & {settling_time_pend:.2f} & {theta_error_fl[-1]:.4f} rad & $k_p = 20, k_d = 8$ & Exponential \\\\")
print(f"Sliding Mode & {reaching_time_smc:.2f} & --- & $k = 5.0, \\lambda = 2.0$ & Finite-time \\\\")
print(f"Backstepping & --- & {final_error_back:.4f} & $c_1 = 2.0, c_2 = 3.0$ & Asymptotic \\\\")

print(r"\bottomrule")
print(r"\end{tabular}")
print(r"\label{tab:performance}")
print(r"\end{table}")
\end{pycode}

\subsection{Equilibrium Point Analysis}

\begin{example}[Equilibrium Classification]
For the controlled Lyapunov system, linearization about $\mathbf{x}_e = \mathbf{0}$ yields Jacobian:
\begin{equation}
J = \begin{bmatrix} -1 & 1 \\ 0 & -1 \end{bmatrix}
\end{equation}
with eigenvalues $\lambda_1 = \lambda_2 = -1$ (repeated real negative), classifying the equilibrium
as a \textbf{stable degenerate node}. The Lyapunov function confirms global asymptotic stability.
\end{example}

\subsection{Control Design Trade-offs}

\begin{pycode}
print(r"\begin{table}[htbp]")
print(r"\centering")
print(r"\caption{Comparative Analysis of Nonlinear Control Techniques}")
print(r"\begin{tabular}{lccccc}")
print(r"\toprule")
print(r"Method & Model Accuracy & Robustness & Chattering & Computational Cost & Applicability \\")
print(r"\midrule")
print(r"Lyapunov Direct & Medium & High & None & Low & Global stability \\")
print(r"Feedback Linearization & High & Low & None & Medium & Exact cancellation \\")
print(r"Sliding Mode & Medium & Very High & Yes & Low & Matched uncertainty \\")
print(r"Backstepping & High & Medium & None & High & Strict-feedback \\")
print(r"\bottomrule")
print(r"\end{tabular}")
print(r"\label{tab:comparison}")
print(r"\end{table}")
\end{pycode}

\section{Discussion}

\begin{remark}[Region of Attraction]
While the quadratic Lyapunov function $V = x_1^2 + x_2^2$ establishes local asymptotic stability,
determining the exact region of attraction requires analysis of sublevel sets $\Omega_c = \{V(\mathbf{x}) \leq c\}$
where $\dot{V} < 0$ holds. For the controlled system, numerical simulation suggests global asymptotic
stability, but formal proof requires invariant set analysis or LaSalle's principle.
\end{remark}

\begin{example}[Zero Dynamics]
In feedback linearization, the internal dynamics (zero dynamics) determine stability when $y \equiv 0$.
For the inverted pendulum with output $y = \theta$, zero dynamics are trivial since relative degree
equals system order, ensuring no instability issues. Systems with unstable zero dynamics are
non-minimum phase and cannot be stabilized by input-output linearization alone.
\end{example}

\subsection{Robustness to Disturbances}

\begin{theorem}[Sliding Mode Robustness]
Sliding mode control achieves invariance to matched disturbances. If $\dot{\mathbf{x}} = \mathbf{f}(\mathbf{x}) + \mathbf{g}(\mathbf{x})(u + d)$
with $|d| < d_{\max}$, choosing control gain $k > d_{\max}$ ensures $s \dot{s} < 0$ despite disturbance $d$.
\end{theorem}

\subsection{Lyapunov Equation for Linear Comparison}

\begin{pycode}
# For comparison, solve continuous Lyapunov equation for linear approximation
# A^T P + P A = -Q
A_linear = np.array([[-1, 1], [0, -1]])
Q = np.eye(2)
from scipy.linalg import solve_continuous_lyapunov
P = solve_continuous_lyapunov(A_linear.T, -Q)
eigenvalues_A = np.linalg.eigvals(A_linear)

print(f"Linear system eigenvalues: $\\lambda_1 = {eigenvalues_A[0]:.2f}$, $\\lambda_2 = {eigenvalues_A[1]:.2f}$")
print(f"\nLyapunov matrix $P$ from algebraic equation:")
print(r"\[")
print(f"P = \\begin{{bmatrix}} {P[0,0]:.2f} & {P[0,1]:.2f} \\\\ {P[1,0]:.2f} & {P[1,1]:.2f} \\end{{bmatrix}}")
print(r"\]")
\end{pycode}

\section{Conclusions}

This computational analysis demonstrates the application of nonlinear control theory:

\begin{enumerate}
\item \textbf{Lyapunov-based control} achieves asymptotic stabilization with settling time
$t_s = \py{f"{settling_time_lyap:.2f}"}$ s and final error $\py{f"{final_error_lyap:.4f}"}$,
confirming the theoretical prediction of exponential convergence along $V$-level curves.

\item \textbf{Feedback linearization} stabilizes the inverted pendulum at $\theta = \pi$ with
settling time $t_s = \py{f"{settling_time_pend:.2f}"}$ s using control gains $k_p = 20$, $k_d = 8$,
achieving exact cancellation of the $\sin(\theta)$ nonlinearity.

\item \textbf{Sliding mode control} drives the sliding surface $s = x_2 + 2x_1$ to zero in
$t_r = \py{f"{reaching_time_smc:.2f}"}$ s with smooth saturation approximation eliminating
chattering while maintaining robustness to matched disturbances.

\item \textbf{Backstepping} design for strict-feedback systems yields asymptotic stability with
final error $\py{f"{final_error_back:.4f}"}$ via recursive CLF construction and virtual control errors.

\item Phase plane analysis reveals equilibrium point classification and limit cycle phenomena
(Van der Pol with $\mu = 2.0$) that linear theory cannot predict.
\end{enumerate}

The computed stability margins and control gains provide quantitative design guidelines for
practical implementation, while Lyapunov analysis offers rigorous stability guarantees beyond
local linearization validity.

\section*{References}

\begin{enumerate}
\item Khalil, H.K. \textit{Nonlinear Systems}, 3rd ed. Prentice Hall, 2002.
\item Slotine, J.-J.E., and Li, W. \textit{Applied Nonlinear Control}. Prentice Hall, 1991.
\item Isidori, A. \textit{Nonlinear Control Systems}, 3rd ed. Springer, 1995.
\item Krstić, M., Kanellakopoulos, I., and Kokotović, P. \textit{Nonlinear and Adaptive Control Design}. Wiley, 1995.
\item Utkin, V., Guldner, J., and Shi, J. \textit{Sliding Mode Control in Electro-Mechanical Systems}, 2nd ed. CRC Press, 2009.
\item Sastry, S. \textit{Nonlinear Systems: Analysis, Stability, and Control}. Springer, 1999.
\item Vidyasagar, M. \textit{Nonlinear Systems Analysis}, 2nd ed. SIAM, 2002.
\item Sepulchre, R., Janković, M., and Kokotović, P. \textit{Constructive Nonlinear Control}. Springer, 1997.
\item Schaft, A.J. van der. \textit{$L_2$-Gain and Passivity Techniques in Nonlinear Control}, 2nd ed. Springer, 2000.
\item Freeman, R.A., and Kokotović, P.V. \textit{Robust Nonlinear Control Design: State-Space and Lyapunov Techniques}. Birkhäuser, 1996.
\item Edwards, C., and Spurgeon, S. \textit{Sliding Mode Control: Theory and Applications}. Taylor \& Francis, 1998.
\item Levine, W.S. (ed.) \textit{The Control Handbook}, 2nd ed. CRC Press, 2011.
\item Astrom, K.J., and Wittenmark, B. \textit{Adaptive Control}, 2nd ed. Dover, 2008.
\item Glad, T., and Ljung, L. \textit{Control Theory: Multivariable and Nonlinear Methods}. Taylor \& Francis, 2000.
\item Sontag, E.D. \textit{Mathematical Control Theory: Deterministic Finite Dimensional Systems}, 2nd ed. Springer, 1998.
\item Byrnes, C.I., and Isidori, A. "Asymptotic stabilization of minimum phase nonlinear systems." \textit{IEEE Trans. Automatic Control}, 36(10):1122-1137, 1991.
\item Emelyanov, S.V. "Binary automatic systems with various forms of pulse modulation." \textit{Proceedings of the USSR Academy of Sciences}, 150:1:106-109, 1963.
\item Lyapunov, A.M. "The general problem of the stability of motion." \textit{Int. J. Control}, 55(3):531-534, 1992 (translation of 1892 original).
\end{enumerate}

\end{document}
